\part{Securing Networks}

Unit 3 moves from a single building to the network that connects it to everything else. Chapter
\ref{ch:3.1} opens with the exploits that operate at the network level itself -- packet sniffing,
spoofing, man-in-the-middle positioning, denial of service -- the attacks that do not care what
application is running because they target the plumbing underneath it. Chapter~\ref{ch:3.2}
narrows the focus to wireless networks specifically, where the attacker does not even need a
cable or a badge, just proximity and the right hardware, which is why Wi-Fi security gets its own
chapter instead of a paragraph inside the last one.

The next two chapters shift from attack to defense. Chapter~\ref{ch:3.3} introduces network
segmentation: the idea that a flat network, where compromising one device gives an attacker a
clear path to every other device, is a design failure, and that VLANs and subnetting exist to
contain a breach rather than just prevent one. Chapter~\ref{ch:3.4} covers firewalls and the
access-control rules that decide what traffic is allowed to cross a segment boundary in the first
place, turning the segmentation from Chapter~\ref{ch:3.3} into something actually enforced.

Chapter~\ref{ch:3.5} closes the unit the same way Unit 2 closed: with detection. Segmentation and
firewalls reduce what can go wrong, but an organization still needs to notice when something does,
and this chapter covers the logging, monitoring, and intrusion-detection practices that make that
possible. By the end of this unit the pattern from Unit 2 -- prevent, control, detect -- should
feel less like a physical-security idea and more like a general security idea that just keeps
reappearing.
